% =============================================================================
%  academic_paper_preamble.tex — Academic LaTeX paper preamble
%
%  Used by: academic_paper_latex_skill.md (working papers, journal
%           submissions, research documents).
%
%  HOW TO USE: copy the contents of this file into the top of the main .tex
%  file you produce. This file is a template, not a runtime include — do not
%  `\input{}` it. It begins with `\documentclass{article}` and must be the
%  literal start of the generated .tex file.
%
%  This is an ACADEMIC paper preamble. The document chrome (title,
%  section headers, links) follows academic conventions — black titles,
%  default-blue links via `blue!60!black`. No FMM branding (no FMM header,
%  no IDB logo footer, no PrimaryBlue title).
%
%  However, Andrea's palette IS defined here for use in figures, TikZ
%  diagrams, and pgfplots charts. Figures often travel between an academic
%  working paper and an FMM policy note as the same chart; they use
%  Andrea's palette in both contexts.
%
%  Colors mirror andreas_palette.md. If you add or change a color, update
%  andreas_palette.md AND every consuming file listed in the palette's
%  "Files that mirror this palette" section, in one atomic commit.
%
%  Last updated: May 28, 2026 — v2.0.0 (v3.2: added the seven new palette colors
%  (DarkYellow, MintGreen, BrightBlue, BrightYellow, BrightGreen, DeepGreen, LabelGray)
%  to the \definecolor block, mirroring andreas_palette.md; available for figures/TikZ.
%  Carries the v2.0.0 floor under the two-tracker convention; kit-wide marker reset and
%  kit package version land at v3.2 cycle close. Prior: v1.0 (May 13, 2026) — extracted
%  from academic_paper_latex_skill.md, Andrea's palette added for figure/TikZ use.)
% =============================================================================

\documentclass[11pt,letterpaper]{article}

% --- Page layout -------------------------------------------------------------
\usepackage[top=1in, bottom=1in, left=1in, right=1in]{geometry}

% --- Spacing -----------------------------------------------------------------
\usepackage{setspace}
\onehalfspacing
\raggedbottom

\usepackage{float}

% --- Encoding and fonts (pdflatex) -------------------------------------------
\usepackage[T1]{fontenc}
\usepackage[utf8]{inputenc}
\usepackage{textcomp}

% --- Math --------------------------------------------------------------------
\usepackage{amssymb,amsmath}

% --- Graphics and figures ----------------------------------------------------
\usepackage{graphicx}
\usepackage{subcaption}

\usepackage[most]{tcolorbox}

% --- Captions ----------------------------------------------------------------
\usepackage[
    justification=centering,
    labelfont=bf,
    labelsep=period,
    font=small,
    skip=6pt
]{caption}

% --- Tables ------------------------------------------------------------------
\usepackage{booktabs}
\usepackage{multirow}
\usepackage{tabularx}
\usepackage{longtable}
\usepackage{array}
\usepackage{calc}

% --- Colors — base xcolor for tables; Andrea's palette for figures/TikZ -----
\usepackage[table,xcdraw]{xcolor}

% Andrea's palette — mirrors andreas_palette.md.
% For use in figures, TikZ, and pgfplots. NOT for document chrome (titles,
% section headers, links) — those follow academic conventions.
\definecolor{PrimaryBlue}{RGB}{25, 110, 140}
\definecolor{SecondaryBlue}{RGB}{94, 144, 168}
\definecolor{LightBlue}{RGB}{203, 216, 223}
\definecolor{DarkBlue}{RGB}{0, 75, 112}
\definecolor{SoftBlue}{RGB}{173, 216, 230}
\definecolor{DarkGreen}{RGB}{76, 114, 52}
\definecolor{LightGreen}{RGB}{180, 200, 150}
\definecolor{WarmCream}{RGB}{255, 243, 214}
\definecolor{DarkYellow}{RGB}{247, 226, 114}
\definecolor{MintGreen}{RGB}{159, 208, 163}
\definecolor{BrightBlue}{RGB}{0, 154, 222}
\definecolor{BrightYellow}{RGB}{255, 218, 0}
\definecolor{BrightGreen}{RGB}{141, 186, 56}
\definecolor{DeepGreen}{RGB}{48, 129, 68}
\definecolor{LabelGray}{RGB}{166, 166, 168}

% --- TikZ and pgfplots -------------------------------------------------------
\usepackage{tikz,pgfplots}
% pgfplotsset is optional — include only if the paper uses pgfplots charts
\pgfplotsset{%
   compat=newest,
   tick label style={font=\footnotesize},
   label style={font=\small},
   legend style={font=\small},
   axis x line = center,
   axis y line = center,
   every axis/.style={pin distance=1ex},
   trim axis left
}

% --- Footnotes ---------------------------------------------------------------
\usepackage[bottom]{footmisc}

% --- Multi-column ------------------------------------------------------------
\usepackage{multicol}

% --- Date formatting ---------------------------------------------------------
\usepackage[en-US]{datetime2}
\DTMnewdatestyle{monthyear}{%
  \renewcommand*{\DTMdisplaydate}[4]{\DTMenglishmonthname{##2} ##1}%
  \renewcommand*{\DTMDisplaydate}{\DTMdisplaydate}%
}

% --- Microtype ---------------------------------------------------------------
\usepackage{microtype}

% --- Line numbers (uncomment for draft review) -------------------------------
\usepackage{lineno}
%\linenumbers

% --- Bibliography ------------------------------------------------------------
\usepackage{natbib}
\bibliographystyle{aer}  % or chicago, plainnat, etc. — confirm with Andrea

% --- Hyperlinks (MUST be last package) ---------------------------------------
\usepackage{hyperref}
\hypersetup{
    colorlinks=true,
    linkcolor=blue!60!black,
    citecolor=blue!60!black,
    urlcolor=blue!60!black,
    pdftitle={Paper Title},
    pdfauthor={Author Names}
}

% =============================================================================
%  CUSTOM COMMANDS
% =============================================================================

\newcommand{\figurenote}[1]{%
    \vspace{4pt}%
    \begin{minipage}{0.95\linewidth}%
        \footnotesize \textit{Note:} #1%
    \end{minipage}%
}

\newcommand{\figuresource}[1]{%
    \vspace{2pt}%
    \begin{minipage}{0.95\linewidth}%
        \footnotesize \textit{Source:} #1%
    \end{minipage}%
}

\newcommand{\tablenote}[1]{%
    \vspace{4pt}%
    \begin{minipage}{0.95\linewidth}%
        \footnotesize \textit{Notes:} #1%
    \end{minipage}%
}

% Prevent overfull lines
\setlength{\emergencystretch}{3em}

